\section{Balanced squarefree blocks and tensor damping}
\label{v2:app:squarefree-damping}

This appendix retains a local interpolation mechanism from the squarefree
construction program.  It supplies exact annihilators and quantitative error
bounds.  These statements alone do not construct a homometric pair: a
subsequent construction must still realize the interpolants as compatible
binary subgroup sections with the required boundary values and nesting.

Let $m=p_1\cdots p_k$ be squarefree and odd.  Choose signs
$\epsilon_i\in\{-1,1\}$ satisfying the population balance
\begin{equation}\label{v2:eq:balanced-singletons}
 \sum_{i=1}^k\epsilon_i(p_i-1)=0.
\end{equation}
On $C_m$, define the radial function $\delta$ to equal $\epsilon_i$ on
the shell of exact element order $p_i$, and zero on every other shell.
Let $\mathcal S=\{p_1,\ldots,p_k\}$ be its order support.  In the divisor
poset of $m$, let $D$ and $U$ denote the indicators of the downward and
upward closures of $\mathcal S$, respectively, and put $P=DU$.
The same letters denote the corresponding radial functions on $C_m$.

\begin{lemma}[A balanced singleton block]\label{v2:lem:balanced-block}
The functions just defined have the following properties.
\begin{enumerate}
\item $\delta$ is the difference of two binary divisor ideals of equal
population, and $P=|\delta|$.
\item The constant function $q=1/2$ satisfies $\delta*q=0$ and
\[
 \|D-q\|_\infty=\|U-q\|_\infty=\|P-q\|_\infty=1/2.
\]
\item The function $\rho=1-D+U$ is increasing in physical order under
divisibility and takes values in $\{0,1,2\}$.  If $d\mid e$ and
$\rho(d)=\rho(e)$, then $D(d)=D(e)$, $U(d)=U(e)$ and $P(d)=P(e)$.
\end{enumerate}
\end{lemma}
\begin{proof}
The positive and negative signs are both present by
\eqref{v2:eq:balanced-singletons}, since every $p_i-1$ is positive.
The two order ideals
\[
 I_+=\{1\}\cup\{p_i:\epsilon_i=1\},\qquad
 I_-=\{1\}\cup\{p_i:\epsilon_i=-1\}
\]
give indicator difference $\delta$.  The shell of order $p_i$ contains
$p_i-1$ elements, so \eqref{v2:eq:balanced-singletons} says
$\sum_x\delta(x)=0$ and $|A_{I_+}|=|A_{I_-}|$.
Since $\mathcal S$ is an antichain, a divisor both above and below
$\mathcal S$ must belong to $\mathcal S$: if $s\mid d\mid t$ with
$s,t\in\mathcal S$, then $s=t=d$.  Therefore $DU=|\delta|$.

For every $x$, $(\delta*q)(x)=\tfrac12\sum_y\delta(y)=0$.
Each of $D,U,P$ is binary, proving the three exact error bounds.
If $d\mid e$, then $D(d)\ge D(e)$ and $U(d)\le U(e)$, and hence
\[
 \rho(e)-\rho(d)=D(d)-D(e)+U(e)-U(d)\ge0.
\]
Equality forces both nonnegative terms to vanish and consequently also
preserves their product.  The possible values of $1-D+U$ are $0,1,2$.
\end{proof}

Examples of \eqref{v2:eq:balanced-singletons} are the signed triples
$+3,+5,-7$ and $+11,+13,-23$, and the signed quartet
$+5,+13,-7,-11$.  For example, the first has populations
$2+4-6=0$.  Equal population is only an ingredient; it is not a claim
of homometry for these two small divisor ideals.  Indeed, in $C_{105}$
the order ideals $\{1,3,5\}$ and $\{1,7\}$ both have population $7$,
but their autocorrelations at the shift $35$ are $3$ and $0$, respectively.

\begin{lemma}[Tensor error cancellation]\label{v2:lem:tensor-damping}
For $1\le i\le t$, let $G_i$ be a finite abelian group, let
$\delta_i:G_i\to\mathbb Q$, and let $b_i,q_i:G_i\to\mathbb Q$ satisfy
$\delta_i*q_i=0$.  On $G=\prod_iG_i$, put
\[
 \Delta(x)=\prod_{i=1}^t\delta_i(x_i),\qquad
 B(x)=\prod_{i=1}^tb_i(x_i),\qquad
 W(x)=B(x)-\prod_{i=1}^t\bigl(b_i(x_i)-q_i(x_i)\bigr).
\]
Then $\Delta*W=0$ and
\begin{equation}\label{v2:eq:damping-error}
 \|W-B\|_\infty\le
       \prod_{i=1}^t\|b_i-q_i\|_\infty.
\end{equation}
In particular, for pairwise coprime balanced singleton blocks and any
independent choices $b_i\in\{D_i,U_i,P_i\}$, the choice $q_i=1/2$ gives
an exact rational annihilator $W$ with
$\|W-B\|_\infty=2^{-t}$.  Its values lie in $2^{-t}\mathbb Z$.
\end{lemma}
\begin{proof}
For tensor functions, finite summation gives
\[
 (\Delta*\textstyle\prod_i f_i)(x)
      =\prod_i(\delta_i*f_i)(x_i).
\]
Expanding the second product in the definition of $W$ cancels the term
$\prod_i b_i$.  Every remaining term contains $q_i$ in at least one
coordinate, so its convolution with $\Delta$ vanishes by the displayed
factorization.  Summing proves $\Delta*W=0$.  The error is exactly
$-\prod_i(b_i-q_i)$, proving \eqref{v2:eq:damping-error}.
For the balanced blocks, every factor $b_i-1/2$ is $1/2$ or $-1/2$ at
every point, so the error has absolute value exactly $2^{-t}$.
Its dyadic denominator and the integer values of $B$ prove the final
claim.  Pairwise coprimality identifies the product of the cyclic blocks
with a cyclic group of squarefree odd order.
\end{proof}

For tensor products of the order antichains, the downward closure,
upward closure and support indicators are respectively
$\prod_iD_i$, $\prod_iU_i$ and $\prod_iP_i$.  Thus the preceding lemma
simultaneously supplies the three approximate boundary indicators used
in interpolation.  A fixed extra factor with error at most $\eta$ gives
error at most $\eta2^{-t}$ by the same lemma, provided its proposed
approximant also satisfies the stated annihilation equation.

\begin{proposition}[Ranks and omitted error factors]
\label{v2:prop:damping-rank}
On a product of $t$ balanced singleton blocks, set
$R=\sum_i\rho_i$.  Then $R$ is increasing, lies in $\{0,\ldots,2t\}$,
and comparable points of equal rank have the same values of every
$D_i,U_i,P_i$.  Consequently every function of this full collection of
indicators, including every $W$ in Lemma~\ref{v2:lem:tensor-damping},
is constant on comparable points of equal rank.

If $S\subseteq\{1,\ldots,t\}$ and $|a_i|\le A_i$ for $i\in S$, then
for arbitrary choices $b_i\in\{D_i,U_i,P_i\}$,
\begin{equation}\label{v2:eq:omitted-factor}
 \left\|\prod_{i\in S}a_i\prod_{i\notin S}(b_i-1/2)\right\|_\infty
 \le \left(\prod_{i\in S}A_i\right)2^{-(t-|S|)}.
\end{equation}
Thus omitting one half-error factor costs exactly a factor of two in the
product bound before accounting for the replacement's size.
\end{proposition}
\begin{proof}
Along a comparable pair each $\rho_i$ has a nonnegative increment.
Equality of their sum forces equality of each increment, so
Lemma~\ref{v2:lem:balanced-block} preserves all the listed indicators.
The assertion about their functions follows immediately.  Each retained
error factor has absolute value $1/2$; multiplication proves
\eqref{v2:eq:omitted-factor}.
\end{proof}

The rank property is a useful compatibility condition when decoding
interpolated sizes into nested sections.  It does not imply that an
arbitrary annihilator $W$ is itself binary, positive, or monotone, nor
that multiplication of $W$ by a rank function preserves annihilation.
Such additional conclusions require separate arguments.  In particular,
the lemmas in this appendix do not certify the larger 196-prime
squarefree construction or its spectator decoder; those require the
additional exact parameters and realization dependencies identified in
the referee guide.
